% basics-ar.tex — Arabic edition of ch. 4, how sourdough works
%
% Original: The Sourdough Framework by Hendrik Kleinwächter,
% the-sourdough-framework.com, CC BY-SA 4.0. The Arabic edition inherits
% CC BY-SA 4.0.
%
% Arabic layer of upstream basics/how-sourdough-works.tex (scaffold
% structure, \label{ch:how-sourdough-works} and
% \label{sec:aggressive-yeast} with its \phantomsection anchor, figure
% labels, \cite keys, footnote positions and the quoting epigraph
% verbatim; text re-typeset from the recovery reference
% text/04-how-sourdough-works.txt, Phase 28R unit 6 — no retranslation).
% Latin terms ride \textenglish{} islands with the parens inside, per
% the preface convention; binomial species names keep the upstream
% \emph{} italics inside their islands. \qty{}/\qtyrange{} ride \LR{}
% islands (the ch. 7 \pHvalue precedent: bare siunitx output mirrors in
% RTL); \num{1500} stands bare (single digit run, no range). The range
% "1--2 ساعة" is wrapped \LR{1--2} per R1-N1. The accepted edition's
% reading "ولسنوات ستتضح قريباً" is kept verbatim (confirmed in the
% accepted PDF text layer, pp. 22–23), notwithstanding the upstream
% "For reasons that will soon be made clear".
%
% Figures: 4 in this chapter, ALL raster (upstream artwork carried over
% unchanged, upstream widths verbatim: 0.7/0.85/0.8/1.0 \textwidth) —
% no TikZ in ch. 4, so the R1-FIGURES route and the ruled uniform-
% Pagella figure digit/Latin face do not arise here.
%
% Short captions: figs 4.1 and 4.3 carry upstream short forms,
% translated; fig 4.4's short form is upstream's own ("Bacteria under
% the microscope"); fig 4.2 has none upstream and gets none. No
% derivations were needed beyond the three upstream short forms.

\chapter{كيف يعمل الخبز الحامض}%
\label{ch:how-sourdough-works}
\begin{quoting}
    نغطي في هذا الفصل أساسيات تخمّر الخبز الحامض. ننظر أولاً إلى التفاعلات الإنزيمية التي تجري في الدقيق لحظة إضافة الماء إليه، وهي التفاعلات التي تستحثّ عملية التخمير. ثم، لفهم هذه العملية فهماً أعمق، نتعرّف أكثر على الكائنات الدقيقة المعنية: الخمائر والبكتيريا.
\end{quoting}

\begin{figure}[!htb]
  \centering
  \includegraphics[width=0.7\textwidth]{infographic-enzymes}
  \caption[تفاعل الأميلازات مع الدقيق]{كيف تتفاعل الأميلازات والبروتييزات مع الدقيق.}%
  \label{fig:infographic-enzymes}
\end{figure}

\section{التفاعلات الإنزيمية}

لفهم التفاعلات الإنزيمية الكثيرة التي تجري عند خلط الدقيق بالماء، علينا أولاً أن نفهم البذور ودورها في دورة حياة القمح وسائر الحبوب.

البذور هي الوسيلة الأساسية التي تتكاثر بها نباتات كثيرة، ومنها القمح. تحوي كل بذرة جنين نباتٍ جديد، ولذا لا بد أن تحوي كل العناصر الغذائية التي سيحتاجها ذلك النبات لينمو.

حين تكون البذرة جافة فهي في وضع السبات، وقد تُخزَّن أحياناً لسنوات عدة. لكن في اللحظة التي تلامس فيها الماء تبدأ بالإنبات. تتحول البذرة إلى جنين نامٍ يحتاج إلى تحويل العناصر المخزونة إلى شيء يستطيع النبات استعماله وهو ينمو. الماء هو ما يستحثّ هذه التفاعلات، والإنزيمات داخل البذرة هي المحفّزات التي تقودها.

تحمل البذرة عادةً أول أوراق النبات الأولية، ويمكنها أن تُرسّخ جذورها مستعينة بالعناصر المخزونة داخلها. وحالما تخترق تلك الأوراق التربة وتلامس ضوء الشمس فوقها تبدأ بالتمثيل الضوئي. هذه العملية هي محرك النبات؛ وبالطاقة التي ينتجها التمثيل الضوئي يستطيع النبات مواصلة نمو جذوره، فيحصل على عناصر إضافية من التربة. وتتيح له تلك العناصر أن ينمو أوراقاً أكثر، فيزداد نشاطه الضوئي حتى يزدهر في بيئته الجديدة.

طبعاً لم يعد الدقيق المطحون قادراً على الإنبات، لكن الإنزيمات التي تستحثّ هذه العملية ما زالت موجودة فيه. ولهذا ينبغي ألا يُطحن الحبوب عند حرارة مرتفعة أكثر من اللازم، لأن ذلك قد يتلف بعض هذه الإنزيمات\footnote{في دراسة حديثة~\cite{milling+commercial+home+mill+comparison} أظهرت الاختبارات أن طحن الدقيق في المنزل بمطحنة صغيرة لم يكن له أثر سلبي يذكر على جودة الخبز الناتج مقارنةً بالدقيق المطحون في مطاحن صناعية كبرى منظّمة الحرارة.}.

عادةً تسدّ بذرة الحبة على الجنين الطريق أمام الممرضات. لكن حين تُطحن الحبة وتتحول إلى دقيق، تنكشف محتويات البذرة. وهذا مثالي لكائناتنا الدقيقة في الخبز الحامض.

فلا الخميرة ولا البكتيريا تصنع غذاءها بنفسها. لكن حين تستحثّ الإنزيمات يصبح الغذاء الذي تحتاجه متوافراً، فتقتات وتتكاثر.

الإنزيمان الرئيسيان في هذه العملية هما الأميلاز \textenglish{(amylase)} والبروتييز \textenglish{(protease)}. ولسنوات ستتضح قريباً، فهما من الأهمية بمكان بالنسبة إلى الخباز المنزلي، ودورهما في صناعة الخبز الحامض قطعة مفتاحية في لغز صنع خبز أطيب طعماً.

\subsection{الأميلاز}

أحياناً، حين تمضغ بطاطا أو قطعة خبز مدةً طويلة، تشعر بطعم حلو على لسانك. السبب أن غدك اللعابي يفرز الأميلاز. يفكك الأميلاز جزيئات النشاء المعقدة إلى سكريات سهلة الهضم. يستعمل جسمك الأميلاز لبدء عملية الهضم، ويعمل الجنين في النبات بالطريقة نفسها مستعملاً الإنزيم عينه: يتخذ من الأميلاز أداةً لصنع سكريات من النشاء، ليصنع منها بدورها مادة نباتية جديدة.

عادةً تعجز الكائنات الدقيقة الساكنة على سطح الحبة عن الإمساك بجزيئات المالتوز المحررة، لأنها تبقى مخبأة داخل الجنين. لكن حين نطحن الدقيق يقع اقتاتٌ محموم. وعموماً كلما ارتفعت الحرارة جرت هذه التفاعلة أسرع. ومع ذلك يحتاج الأميلاز إلى وقت ليحوّل معظم النشاء إلى سكريات بسيطة، وهي سكريات لا تلتهمها الخميرة فحسب، بل هي ضرورية أيضاً لتفاعل مايار \textenglish{(Maillard reaction)} المسؤول عن الاسمرار المحسّن أثناء الخَبز. ولهذا كان التخمير الطويل مفتاح صناعة خبز عظيم.

إن كنت من هواة تخمير الجِعة فتعلم أن من المهم إبقاء الجِعة عند درجات حرارة محددة ليتحول النشاء فيها إلى سكر بفعل الأميلازات المختلفة~\cite{beer+amylase}. هذه العملية مهمة إلى حدّ أن هناك اختباراً شائعاً يُعرف به هل تحوّل النشاء كله أم لا. يسمى هذا الاختبار اختبار النشاء باليود \textenglish{(Iodine Starch Test)}، ويقتضي خلط اليود في عينة من خلطتك ومراقبة اللون: فإن كان أزرق أو أسود فما زال فيها نشاء لم يتحول. ويتردد في ذهني: هل ينفع اختبار مماثل مع عجين الخبز؟

يواجه الخبازون الصناعيون، الذين يضيفون خميرة فائقة النشاط لإنتاج الخبز في زمن قصير، مشكلة مشابهة. حلهم أن يضيفوا دقيقاً مُنبتاً إلى العجين؛ وهذا الدقيق غني بالإنزيمات فيسرّع بذلك عملية التخمير. في المرة القادمة التي تتسوق فيها، فتّش غلاف الخبز الذي تشتريه: إن وجدت كلمة «مَلْت» \textenglish{(malt)} في قائمة المكوّنات، فالمرجح أن هذه الاستراتيجية كانت مستعملة.

لاحظ أن للملت صنفين فعلاً: صنف مَلْت إنزيميّ النشاط \textenglish{(enzymatically active malt)} لم يُسخَّن فوق \LR{\qty{70}{\degreeCelsius}}، وهي الحرارة التي يبدأ عندها الأميلاز بالتدهور؛ وصنف مَلْت خامد \textenglish{(inactive malt)} سُخّن إلى حرارة أعلى، فلا أثر له على دقيقك.

\subsection{البروتييز}

كما يفكك الأميلاز النشاء إلى سكريات بسيطة، يفكك البروتييز البروتينات المعقدة إلى بروتينات أبسط وأحماض أمينية. ولأن القمح يحوي الغلوتين، وهو بروتين أساسي لبنية الخبز، فلا بد أن يؤدي البروتييز دوراً محورياً في خَبز الخبز الحامض.

ولأن بذور الحبوب تحتاج إلى أحماض أمينية أصغر لبناء الجذور وسائر مواد النبات، يبدأ الغلوتين في تلك البذور بالتفتت لحظة إنباتها؛ ولما كان إضافة الماء إلى الدقيق تستحثّ الإنزيمات عينها، تجري العملية نفسها في عجين الخبز.

إن حاولت يوماً صنع عجين قمحي وإبقائه في حرارة الغرفة عدة أيام، فستكتشف بنفسك أن شبكة الغلوتين تتفكك حتى يعجز العجين عن التماسك. عندها يتمزق العجين بسهولة، ولا يحفظ بنية، ولم يعد صالحاً لخبز الخبز.

حدث لي هذا ذات مرة حين حاولت صنع بادئة خبز حامض مباشرةً من بادئة مجففة. في اليوم الثالث إلى الرابع كان التخمير بطيئاً لدرجة أن شبكة الغلوتين تفككت. وكان السبب الجذري لهذه المشكلة هو البروتييز: فما تضيفه من ماء إلى عجينك يستحثّ البروتييز، فيشتغل في تهيئة الأحماض الأمينية للجنين.

وهذه تجربة ممتعة أخرى يمكنك تجربتها لتصوّر أهمية البروتييز: اصنع عجيناً سريع التخمير بكمية كبيرة من الخميرة الجافة النشطة. خلال \LR{1--2} ساعة ينتفخ عجينك ويزداد حجمه. اخبزه ثم افحص بنية اللُّب: ستجده كثيفاً تماماً، بعيداً بكل بعيد عن ما يمكن أن يبلغه من رغوة. السبب أن إنزيم البروتييز لم يُمنح وقتاً كافياً لأداء عمله.

في البدء، أثناء العجن، يصير العجين مرناً ويتماسك جيداً. لكنه مع التخمير يصير أرخى وأكثر قابلية للمد~\cite{protease+enzyme+bread}. والسبب أن بعض روابط الغلوتين تكون قد تفككت طبيعياً بفعل البروتييز في عملية تُعرف بالتحلل البروتيني \textenglish{(proteolysis)}. هذا ما ييسّر على الخميرة نفخ العجين، ولهذا كان التخمير الطويل حاسماً حين تريد لُباً رغوياً مفتوح المسامات في خبزك الحامض.

وإلى جانب استعمال مكوّنات ممتازة، فإن التخمير البطيء من الأسباب الرئيسية لما يتذوقه الناس في البيتزا النابولية: فالبروتييز يجعل العجين قابلاً للمد وسهل النفخ، فيتحقق حافة طرية هوائية.

ولأن عملية التخمير تأخذ عادةً أكثر من 8 ساعات، ينبغي استعمال دقيق أعلى محتوى بالغلوتين. فهذا يمنح العجين وقتاً أطول ليتفكك بفعل البروتييز دون أن يمس مرونته. وإن استعملت دقيقاً أضعف فقد ينتهي بك الأمر إلى عجين تفكك كثيراً حتى يتمزق عند المد، فيستحيل مثلاً تشكيله إلى قرص بيتزا.

كانت البيتزا تقليدياً تُصنع بالخبز الحامض، أما حديثاً فتُصنع بالخميرة الجافة النشطة؛ إذ بقاء العجين صالحاً مدةً أطول يجعل التعامل معه أسهل على المستوى التجاري. وإن استعملت الخبز الحامض فلك أن تنال نافذة من ثلاثين إلى تسعين دقيقة قبل أن يبدأ العجين بالتدهور، بسبب اشتغال البروتييز مدةً أطول، وبسبب نواتج البكتيريا الثانوية التي سنتناولها بتفصيل أوضح لاحقاً في هذا الفصل.

\subsection{تحسين النشاط الإنزيمي}

كما سبق شرحه، المَلْت حيلة شائعة لتسريع النشاط الإنزيمي. لكنني شخصياً أفضّل تجنّب المَلْت واستعمال حيلة تعلمتها وأنا أصنع خبز القمح الكامل.

حين بدأت أول مرة بصنع خبز القمح الكامل، لم أكن أبلغ القشرة ولا اللُّب ولا القوام الذين أريدهما مهما جرّبت. كان عجيني يميل بدلاً من ذلك إلى فرط التخمير بسرعة. أما عند استعمال دقيق أبيض بمحتوى غلوتين مشابه، كان خبزي يخرج عظيماً دائماً.

كنت آنذاك أستعمل أوتولايزاً ممتداً، وهي كلمة فخمة لمعنى بسيط: خلط الدقيق والماء مسبقاً ثم ترك الخليط مرتاحاً. تستدعي معظم الوصفات هذه الخطوة لأنها تمنح العجين بدايةً إنزيمية مبكرة، وهي فكرة عظيمة عموماً. لكن بديلاً مماثل الفاعلية يمكنك ببساطة تخفيض كمية عامل التخمير المستعمل، أي بادئتك في حالة الخبز الحامض. يتيح هذا للتفاعلات الحيوية نفسها أن تجري بالسرعة نفسها تقريباً دون أن تضطر إلى خلط عجينك عدة مرات. تحسّن مستوى خبز القمح الكامل عندي تحسناً هائلاً بعد أن توقفت عن إخضاع عجائني للأوتولايز.

والآن بعد أن أجدت التفكير في الأمر، النتيجة التي لاحظتها تبدو منطقية. في الطبيعة تلامس الأجزاء الخارجية للبذرة الماء أولاً، ولا يجد الماء طريقه ببطء إلى مركز الحبة إلا بعد أن يخترق هذا الحاجز. تحتاج البذرة أن تنبت أولاً كي تتفوق على البذور المجاورة لها، وهذا يستلزم دخول الماء بسرعة. لكن عليها في الوقت نفسه أن تدافع عن نفسها ضد الحيوانات، وضد البكتيريا والفطريات الخطر المحتمل، وهذا يستلزم حاجزاً يحمي الجنين الداخلي. وسيلة النبات للتوفيق بين الهدفين هي أن يوجد معظم الإنزيمات في الأجزاء الخارجية من القشرة، فهي تستحثّ أولاً~\cite{enzymatic+activity+whole+wheat}. لهذا، بمجرد أن تضيف قليلاً من دقيق القمح الكامل إلى عجينك، ينبغي أن تتمكن من تحسين نشاطه الإنزيمي تحسناً واضحاً. ولهذا أضيف عادةً إلى عجائن الدقيق الأبيض العادي \LR{\qtyrange{10}{20}{\percent}} من دقيق القمح الكامل.

\begin{figure}
  \centering
  \includegraphics[width=0.85\textwidth]{whole-wheat-crumb}
  \caption{خبز حامض من القمح الكامل.}%
  \label{fig:whole-wheat-crumb}
\end{figure}

بفهمك للإنزيمين الرئيسيين، الأميلاز والبروتييز، تصبح أفضل تجهيزاً لصنع خبز على ذوقك. تفضّل لُباً أطرى أم أصلب؟ تريد قشرة أفتح أم أغمق؟ ترغب في تخفيض كمية الغلوتين في خبزك النهائي؟ كلها عوامل يمكنك ضبطها بمجرد تعديل سرعة تخمير عجينك.

\section{الخمائر}

الخمائر كائنات دقيقة وحيدة الخلية تنتمي إلى مملكة الفطريات. تتكاثر بالتنبّس أو ببناء الأبواغ. الأبواغ متناهية في الصغر وشديدة المقاومة للعوامل الخارجية؛ فوجد العلماء أبواغاً سليمة عمرها مئات ملايين السنين. وتتنوع أنواعها تنوعاً واسعاً: حتى الآن حُدد نحو \num{1500} نوع. وعلى خلاف سائر أفراد مملكة الفطريات كالعفن، لا تبني الخمائر عادةً شبكة غزل فطرية~\cite{molecular+mechanisms+yeast}.\footnote{عن استثناء مثير لاهتمام، انتقل إلى آخر هذا القسم في صفحة~\pageref{sec:aggressive-yeast}.}

\begin{figure}[!htb]
  \centering
  \includegraphics[width=0.8\textwidth]{saccharomyces-cerevisiae-microscope}
  \caption[الخميرة الجِعية تحت المجهر]{الخميرة الجِعية \textenglish{Saccharomyces cerevisiae} تحت المجهر.}%
  \label{fig:saccharomyces-cerevisiae-microscope}
\end{figure}

الخمائر فطريات رمية التغذية \textenglish{(saprotrophic)}، أي أنها لا تصنع غذاءها بنفسها، بل تعتمد على مصادر خارجية تفككها وتحللها إلى مركبات يسهل أيضها. ما نسميه اليوم عملية التخمير هو تحلل الخميرة للكربوهيدرات إلى ثاني أكسيد الكربون وكحول. هذه العملية معروفة منذ آلاف السنين، واستعملت منذ القِدم في صناعة الخبز وكذلك المشروبات الكحولية.

تستطيع الخميرة النمو والتكاثر في شروط هوائية ولا هوائية على السواء. حين يتوفر الأكسجين تنتج ثاني أكسيد الكربون والماء حصراً تقريباً، وحين يغيب يتحول أيضها إلى إنتاج مركبات كحولية~\cite{effects+oxygen+yeast+growth}.

تتفاوت درجات الحرارة التي تنمو عندها الخمائر. بعض الخمائر، مثل \textenglish{\emph{Leucosporidium frigidum}}، تنمو أفضل ما يكون بين \LR{\qty{-2}{\degreeCelsius}} و\LR{\qty{20}{\degreeCelsius}}، بينما تفضّل أخرى حرارة أعلى. وعموماً كلما دفأ الوسط أسرع أيض الخميرة. صنف الخميرة الذي تربّيه في بادئتك ينبغي أن يعمل أفضل ما يكون ضمن مدى الحرارة التي نبتت الحبة وحُصدت فيها. فإذا كنت من مكان بارد وربّيت بادئة من صنف جاودار إسكندنافي، فالمرجح أن تفضّل خميرتك وسطاً أبرد.

مثالاً على ذلك: اكتشف صانعو الجِعة خميرة نافعة تعيش في الكهوف الباردة حول مدينة بلزن في جمهورية التشيك. صارت هذه الخميرة معروفة منذ ذلك الحين بإنتاجها جِعة ممتازة عند حرارة منخفضة، وتستعمل أصناف سلالاتها اليوم في تخمير أنواع اللاغر الشهيرة.

الخمائر عموماً كائنات شائعة جداً. توجد على حبوب الحبوب والفواكه وكثير غيرها من نباتات التربة، بل توجد داخل أمعائك! وصدف أن الخمائر التي نستعملها للخبز تُربّى على أوراق النباتات، مع أن معرفتنا بما بها من بيئة لا تزال شحيحة.

تحتمي النباتات بجدران خلوية سميكة يتعذر على قلة من الفطريات أو البكتيريا اختراقها. لكن بعض الأنواع تنتج إنزيمات قادرة على تفكيك تلك الجدران كي تصيب النبات.

بعض الفطريات والبكتيريا تعيش داخل النباتات دون أن تسبب لها أي أذى. هذه تُعرف بالكائنات الداخلية \textenglish{(endophytes)}. فهي لا تضرّ مضيفها، بل تعيش معه في علاقة تعايش: تعين النباتات التي تسكنها على التحصّن من ممرضات أخرى قد تصيبها عبر أوراقها. وإلى جانب هذه الحماية تعينها أيضاً على تحمّل الإجهاد المائي والحراري، وعلى توافر العناصر الغذائية. ومقابل خدمتها لنباتاتها المضيفة تنال هذه الفطريات والبكتيريا كربوناً للطاقة.

لكن العلاقة بين الداخلية والنبات ليست دائماً متبادلة النفع: فأحياناً، تحت الإجهاد، تتحول إلى ممرض غازي تفني مضيفه بالتحلل في نهاية المطاف~\cite{endophytes+in+plants}.

وهناك كائنات دقيقة أخرى، على خلاف الداخلية، لا تخترق الجدران الخلوية بل تعيش على سطح النبات وتنال غذاءها من مياه المطر أو الهواء أو حيوانات أخرى. بعضها يقتات الندى العسلي الذي يفرزه المنّ، أو حبوب الطلع التي تهبط على سطح الأوراق. تسمى هذه الكائنات النابِرات \textenglish{(epiphytes)}، ومنها الخمائر التي نستعملها للخبز.

واللافت أنه إذا أُزيلت المصادر الغذائية الخارجية ما زال عدد كبير من الفطريات والبكتيريا النابرة يُعثر عليه على سطح النبات، ما يوحي بأنها لا بد أنها تقتات من النبات مباشرة. وفعلاً تشير بعض الأبحاث إلى أن بعض النباتات تُطلق عمداً على سطحها مركبات من سكريات وأحماض عضوية وأمينية وميثانول وأملاح متنوعة، فتجذب هذه العناصر النابرات الساكنة على سطح النبات.

النابرات نافعة لبقاء النبات، إذ تمنحه حمايةً محسّنة ضد العفن وسائر الممرضات. وفعلاً من مصلحة النابرات أن تبقي نباتاتها المضيفة حية أطول ما يمكن~\cite{leaf+surface+sugars+epiphytes}.

يُجرى كل يوم بحثٌ أكثر في سبل استعمال الخمائر عواملَ مكافحةٍ حيوية \textenglish{(biocontrol agents)} لحماية النبات، وميزتها أن هذه العوامل الحيوية آمنة غذائياً، لأن سلالات الخميرة المعنية تُعدّ عموماً غير ضارة بالناس. ستنمو الخمائر وتتكاثر على الأوراق فتحجبها في جوهرها عن سائر أصناف العفن. وقد يكون هذا نقلة تغيّر قواعد اللعبة للمزارع الكرومية التي يرهقها البياض الدقيقي.

ويمكن استعمال هذه العوامل الحيوية أيضاً لتحصين النباتات ضد فطر الأرغوت \textenglish{(ergot)} المؤثر نفسياً، وهو فطر يفضّل النمو في الأوساط الأبرد والأكثر رطوبة، ويشكّل مشكلة بالغة لمزارعي الجاودار. خفّض المشرّعون مؤخراً الحد المسموح به من تلوث الأرغوت في دقيق الجاودار، لأنه يصيب الحبة ويجعلها غير صالحة للأكل بسبب سُمّيته العالية للكبد. ويمكن للخمائر أن تعين على تخفيف تلوث الأرغوت.

وهناك تجربة أخرى مثيرة أجراها علماء إيطاليون تبيّن ما لِلخمائر من قدر في حماية محاصيلنا. أحدثوا أولاً شقوقاً دقيقة في بعض عنبات العنب على الكرمة، ثم أصابوا الجروح بالعفن. شقوق لم تصب إلا بالعفن، وأخرى لُقّحت إضافةً إلى العفن ببعض سلالات الخمائر البرية البالغ عددها 150 المعزولة من الأوراق. وجدوا أنه حين لُقحت الجروح بالخميرة لم يصب العنب بأذى يُذكر~\cite{yeasts+biocontrol+agent}.

\phantomsection~\label{sec:aggressive-yeast}%
والأعجب أن تجربة أخرى أظهرت أن خميرة الجِعة قد تعمل ممرضاً عدوانياً على الكرمات. بدأت الخميرة حياتها في تعايش مع النباتات، لكن بعد أن أصيبت الكرمات بأضرار بالغة صارت الخميرة انتهازية بدأت تهاجم، حتى إنها أنتجت خيوطاً فطرية، هي شبكة الغزل الفطرية المألوفة عند الفطريات، لتتمكن من اختراق أنسجة النبات.

\section{البكتيريا}

الخصم الميكروبي الآخر الأكثر هيمنة في خبزك الحامض هو البكتيريا. بل إنها مهيمنة إلى حد أنها تفوق الخميرة في عجينك بنسبة 100 إلى 1. فحيث توفر الخميرة قوة النفخ، تصنع البكتيريا النكهات المميزة التي سُمي الخبز الحامض باسمها. وتعود هذه النكهات إلى النواتج الحمضية الثانوية الناتجة عن اقتات البكتيريا. ومن مزايا هذه الحموض أنها تمدّ مدة صلاحية الخبز الحامض تمدّداً كبيراً~\cite{shelflife+acidity}.

\begin{figure}
  \centering
  \includegraphics[width=1.0\textwidth]{bacteria-microscope}
  \caption[البكتيريا تحت المجهر]{\textenglish{Fructilactobacillus sanfranciscensis} تحت المجهر.}%
  \label{fig:lactobacillus-franciscensis-microscope}
\end{figure}

يوجد صنفان سائدان من الحموض في الخبز الحامض: حمض اللبنيك وحمض الخليك. أما في النكهة فللحمض اللبنيك نوتات ألبان واضحة، بينما حمض الخليك طعمه طعم الخل (وهو فعلاً مكوّنه الرئيسي!). وتنتج هذه النواتج الحمضية بكتيريا حمض اللبنيك متجانسة التخمير \textenglish{(homofermentative)} وغير متجانسته \textenglish{(heterofermentative)}.

يعني «متجانسة التخمير» أن البكتيريا تنتج أثناء التخمير مركباً واحداً، وهو هنا حمض اللبنيك. أما «غير المتجانسة» فتعني أن مركبات أخرى تنتج أيضاً: هنا لا يُنتج حمض اللبنيك فقط، بل حمض الخليك أيضاً، وكذلك الإيثانول بل وقليل من ثاني أكسيد الكربون، وهما ناتجان يُنسبان عادةً إلى الخميرة. ومن أشهر سلالات بكتيريا حمض اللبنيك \textenglish{\emph{Fructilactobacillus sanfranciscensis}}، التي اشتقت اسمها من الخبز الحامض السان فرانسيسكوني الشهير؛ فأول زرعة معزولة منها جاءت من مخبز في هذه المدينة، ومن هنا الاسم.

تتنافس الخميرة والبكتيريا على المصدر الغذائي نفسه: السكر. أفاد بعض العلماء أن البكتيريا تستهلك المالتوز أساساً بينما تفضّل الخميرة الغلوكوز، وأفاد آخرون أن البكتيريا تقتات من نواتج الخميرة الثانوية والعكس صحيح. وهذا منطقي، إذ تقوم الطبيعة عموماً بعمل بارع في التسميد وتفكيك المادة الحيوية~\cite{lactobacillus+sanfrancisco}.

لم أعثر بعد على مصدر سليم يصف بدقة التعايش بين الخميرة والبكتيريا، لكن فهمي الحالي أنهما تتعايشان وتنفع إحداهما الأخرى أحياناً دون غيرها. فالخميرة مثلاً تحتمل الوسط الحمضي الذي تصنعه البكتيريا المحيطة، فتتحصّن بذلك من ممرضات أخرى. لكن في المقابل تُظهر أبحاثٌ أخرى أن كلا الصنفين ينتج مركبات تمنع الآخر من أيض غذائه، وهي ملاحظة مثيرة للاهتمام بالمناسبة، إذ قد تعين على اكتشاف مضادات حيوية أو مبيدات فطرية إضافية~\cite{mold+lactic+acid+bacteria}.

حاولت يوماً تربية الفطر الصالح للأكل ولاحظت الغزل الفطري وهو يدافع عن نفسه ضد البكتيريا المحيطة؛ إذ أنتج كلا الصنفين بفاعليةٍ مركبات لمحاربة الآخر. ومع ذلك بدا القتالُ بعد حين أنه بلغ طريقاً مسدوداً، إذ قد نما الغزل الفطري حول بقعة البكتيريا كاملاً فمنعها من الانتشار أبعد. أتخيّل أن مشهداً مشابهاً قد يجري في بادئات الخبز الحامض لدينا، وإن كان في بيئة الخبز الحامض، وهي أكثر سيولة، لا بد أن يقع هذا القتال في كل مكان من العجين لا في بقعة معزولة. ما زال البحث مطلوباً في هذا الموضوع لفهمٍ أفضل لتفاصيل العلاقة بين الخميرة والبكتيريا.

من السمات الأخرى المثيرة للاهتمام لدى بكتيريا الخبز الحامض قدرتها على تفكيك بروتينات عجينك واستهلاكها. إن كنت قد خبزت الخبز الحامض من قبل فالمرجح أنك جرّبت هذا بنفسك. وتتذكر من قسم «التفاعلات الإنزيمية» أن البروتييز يفكك شبكة الغلوتين في عجينك، فتكون النتيجة فوضى لزجة إذا أُبقي العجين دون خبز مدةً طويلة. والبكتيريا أيضاً تستهلك الغلوتين في عجينك وتفككه في عملية تسمى التحلل البروتيني \textenglish{(proteolysis)}.

كان هذا بالنسبة إليّ لغزاً عظيماً حين بدأت العمل بالخبز الحامض: فهو من جهة يجعل العجين ألزق، ومن جهة أخرى يجعله أكثر قابلية للمد وأسهل تعاملاً. فمع تناقص الغلوتين يصير العجين أسهل على الكائنات الدقيقة أن تنفخه، فيتماسك ويرتفع. يمكن تشبيه ذلك بفرق الجهد بين نفخ إطار مطاطي سميك وبالون رقيق هشّ: الأخير ينفخ بفمك بسهولة، أما الأول فلا.

ولا غرابة أن يُسرّع التحللَ البروتيني إنزيمُ البروتييز الذي سبق الحديث عنه، إذ يعين على تفكيك الغلوتين إلى أحماض أمينية أصغر يسهل أيضها.

هذا بالنسبة إليّ هو عجيب عملية التخمير. حين تأكل خبزاً حامضاً فأنت لا تستهلك دقيقاً وماءً فحسب، بل النتيجة النهائية لعمليات حيوية معقدة أنجزتها البكتيريا والخميرة. ولسداد المكوّن البكتيري الإضافي، يحوي الخبز الحامض عادةً غلوتيناً أقل من عجين يعتمد على الخميرة وحدها~\cite{proteolysis+sourdough+bacteria}. وعلاوة على ذلك، تمثّل البكتيريا متجانسة التخمير الإيثانولَ الذي تنتجه الخميرة وسائر بكتيريا حمض اللبنيك غير المتجانسة، وفي الحالين تكون معظم المركبات الناتجة أحماضاً عضوية. فكل مورد طبيعي في خبزك الحامض تعيد الكائنات الدقيقة داخله تدويره، وهي جميعاً تحاول أكل ما يجد متاحاً أطول ما يمكن، ومع كل تغذية تصبح أكثر براعة في استغلال هذه الموارد.

حسب ملف النكهة الذي تفضّله يمكنك الانتقاء لمصلحة حمضٍ عضوي أو آخر. إنتاج حمض الخليك يحتاج أكسجيناً، فبتحريم بادئتك منه يمكنك تعزيز جمهرة بكتيريا حمض اللبنيك متجانسة التخمير. ومع الزمن ستهيمن تلك البكتيريا وتتفوق على المنتجة لحمض الخليك~\cite{acetic+acid+oxygen}.

درجة حرارة التخمير المثلى لبكتيريا حمض اللبنيك لديك تعتمد على السلالات التي ربّيتها في بادئتك. وعموماً تعمل تلك البكتيريا أفضل ما يكون عند الحرارة التي أُعدت عندها البادئة، لأنك انتقيت أصلاً بالتربية بكتيريا تزدهر في ذلك الشرط.

وفي تجربة جديرة بالذكر فحص العلماء بكتيريا حمض اللبنيك على أوراق الذرة. خفضوا حرارة الوسط من \LR{\qtyrange{20}{25}{\degreeCelsius}} إلى نحو \LR{\qtyrange{5}{10}{\degreeCelsius}}، ثم لاحظوا بعدها أصنافاً من البكتيريا لم يسبق للناس أن رأوا مثلها~\cite{temperature+bacteria+corn}، فيما يؤكد أن تنوعاً كبيراً من سلالات البكتيريا يعيش فعلاً على أوراق النبات.

وبالمناسبة، يمكنك أن تجري تجربة مشابهة ببدء بادئة خبز حامض عند حرارة أدنى. نظرياً ينبغي أن يتكيف الحيز الميكروبي، إذ ستبدأ الكائنات الدقيقة التي تزدهر أكثر من غيرها عند الحرارة المنخفضة بالهيمنة. وسيكون ممتعاً أن نرى هل يمكن لهذا أن يؤثر فعلاً في طعم الخبز الناتج.

ملاحظة أخيرة جديرة بالذكر: تقول بعض المصادر إن التخمير عند حرارة أدنى قد يرفع إنتاج حمض الخليك، بينما إن التخمير عند حرارة أدفأ قد يعزز إنتاج حمض اللبنيك. لم أستطع التحقق من ذلك في تجاربي الخاصة. ما زال البحث مطلوباً في هذا الموضوع.
